\chapter{解題前的準備}
\label{ch:prep}

一間小診所早上收到一疊掛號單，順序是亂的。 護理師有兩種叫號的做法。 第一種：
每次要叫下一位，就把還沒叫的單子從頭翻一遍，找出號碼最小的那張。 第二種：開診
前先把每張單子插進一個照號碼分好格的架子，之後每次叫號只要從架子上依序拿下一
張。 插一張單子比翻一張慢，大約要花翻十張的時間。 病人只有五位時，哪一種比較省
事？ 如果有一百位、一千位呢？ 往下讀之前，先自己估估看，人數到多少時第二種開始
划算。

這個問題沒有用到任何演算法，只用到一件事：數清楚一個做法要做幾次動作，再看這個
次數隨人數怎麼長大。 寫解題程式也是一樣。 動手之前先算程式大約要做多少次運算，
就知道它在題目的時間限制內跑不跑得完；這比寫完、送出、等結果，再回頭改快得多。

這一章不教特定的演算法，教的是後面每一章都會用到的四件事：

\begin{itemize}
\item 數清楚一段程式執行幾次，以及這個次數隨資料量怎麼長大（第~\ref{sec:count-loops}、\ref{sec:growth}~節）。
\item 從題目給的資料範圍，判斷需要多快的做法（第~\ref{sec:growth}~節）。
\item 避開兩種和演算法無關、卻很常讓答案錯掉的問題：整數溢位和輸入輸出（第~\ref{sec:overflow}、\ref{sec:io}~節）。
\item 自己產生測試資料，用一個慢但簡單的程式來檢查快的程式（第~\ref{sec:testing}~節）。
\end{itemize}

這一章會用到《C++ 程式設計：給學過 C 的人》（以下簡稱銜接書）裡的幾個地方：
解題程式的骨架、從資料範圍估計時間的表、用檔案測試、\texttt{long long} 的溢位在
第~10~章 10.1 節；輸入加速的兩行在第~1~章；\texttt{getline} 在第~3~章；
\texttt{vector}、\texttt{pair}、\texttt{sort}、\texttt{set}／\texttt{map}、
\texttt{lower\_bound} 分別在第~3、8、9~章。 本書從第~2~章開始會一直用到這些容器
和函式，章末的實作練習也放了一組複習題。 沒讀過銜接書的讀者，只要會寫一支從標準
輸入讀資料、把答案印到標準輸出的 C++ 程式，就能讀這一章。

\section{數一數程式做了幾次}
\label{sec:count-loops}

先把章首的第一種叫號方法寫成程式。 \texttt{num} 存著 $n$ 張單子的號碼，每次叫號，
都在還沒叫過的單子裡找最小的，找到後換到前面，表示已經叫過：

\begin{code}[label=code:call-loop]{cpp}{第一種叫號方法：每次翻遍還沒叫的單子}
long long flips = 0;                      // 比較了幾次號碼
for (int call = 0; call < n; call++) {    // 第 call 次叫號
    int best = call;
    for (int i = call + 1; i < n; i++) {  // 只翻還沒叫過的
        flips++;
        if (num[i] < num[best]) {
            best = i;
        }
    }
    swap(num[call], num[best]);           // 叫過的換到前面
}
\end{code}

內層迴圈的範圍從 \texttt{call + 1} 開始，跟著外層變。 這種迴圈不能用「外層 $n$ 次
乘內層 $n$ 次」來算，要一輪一輪列出來。 $n = 5$ 時：

\begin{center}
\begin{tabular}{cccc}
\toprule
\texttt{call} & 內層的 \texttt{i} & 這一輪的 \texttt{flips++} & 累計 \\
\midrule
0 & 1, 2, 3, 4 & 4 & 4 \\
1 & 2, 3, 4 & 3 & 7 \\
2 & 3, 4 & 2 & 9 \\
3 & 4 & 1 & 10 \\
4 & （不執行） & 0 & 10 \\
\bottomrule
\end{tabular}
\end{center}

每一輪少一次，加起來是 $4 + 3 + 2 + 1 + 0 = 10$。 一般的 $n$ 是
\[
(n - 1) + (n - 2) + \cdots + 1 + 0 = \frac{n(n - 1)}{2}.
\]
這裡數的是 \texttt{flips++} 這一行，也就是比較號碼的次數。 迴圈本身的 \texttt{i++}、
\texttt{i < n} 也要花時間，但它們的次數跟 \texttt{flips++} 差不多，數最內層那一行
就足夠代表整段程式的工作量。

\subsection{三種常見的迴圈形狀}

大部分迴圈的次數，可以拆成下面三種形狀來算。

\begin{itemize}
\item \textbf{固定範圍。} \texttt{for (i = 0; i < n; i++)} 跑 $n$ 次；兩層都是固定範
圍時相乘，跑 $n \times m$ 次。
\item \textbf{內層範圍跟著外層變。} 像上面的叫號程式，要一輪一輪加起來。 常見的結果
是 $1 + 2 + \cdots + n = n(n + 1)/2$，大約是 $n^2/2$。
\item \textbf{每次乘或除一個數。} \texttt{for (i = 1; i <= n; i *= 2)} 的 \texttt{i}
依序是 $1, 2, 4, 8, \ldots$，到超過 $n$ 為止。 $n = 20$ 時是 $1, 2, 4, 8, 16$，跑
$5$ 次；$n = 1000$ 時只跑 $10$ 次。 一般來說是 $\lfloor \log_2 n \rfloor + 1$
次，也就是 $n$ 寫成二進位的位數。 每次除以 $3$ 的迴圈也一樣，次數大約是
$\log_3 n$。
\end{itemize}

第三種形狀的次數長得非常慢：$n$ 從一千變成一百萬，次數只從 $10$ 變成 $20$。 所以
看到「兩層迴圈」先別急著說它很慢，要看清楚每一層的變數是怎麼變的。 下面的例題就
是三段看起來都是兩層、實際次數卻差很多的程式。

\subsection{例題：迴圈到底跑幾次}
\label{sec:ex-loop-count}

\textbf{題意}（\akiref{how-many-times-does-the-loop-run}{迴圈到底跑幾次}）

\begin{itemize}
\item 有三段程式片段，變數都是 \texttt{long long}，\texttt{cnt} 從 $0$ 開始（見下面的
片語）。
\item 有 $q$ 個問題，每個問題給片段編號 $t$ 和正整數 $n$，問那一段跑完後
\texttt{cnt} 是多少。
\item $q \le 10^5$，$n \le 10^{18}$。 題目保證片段裡的變數不會溢位。
\end{itemize}

\begin{code}[label=code:three-pieces]{cpp}{例題的三段片段}
// 片段 1
for (i = 1; i <= n; i++) {
    j = i;
    while (j <= n) {
        cnt++;
        j *= 2;
    }
}
// 片段 2
for (i = 1; i <= n; i *= 2)
    for (j = 1; j <= i; j++)
        cnt++;
// 片段 3（i /= 3 是整數除法）
for (i = n; i > 0; i /= 3)
    for (j = 0; j <= i; j++)
        cnt++;
\end{code}

例如 $n = 9$ 時，片段 3 的外層 \texttt{i} 依序是 $9, 3, 1$，下一個是 $0$ 就停了。
內層的 \texttt{j} 從 $0$ 跑到 \texttt{i}，三輪分別跑 $10$、$4$、$2$ 次，答案是
$16$。

\textbf{想法}

最直接的做法，是把三段片段原封不動放進程式裡跑：寫一個函式
\texttt{run(t, n)}，裡面用 \texttt{if} 分成三段，每一段就是上面的片語照抄，最後回傳
\texttt{cnt}；\texttt{main} 讀每個問題、印出 \texttt{run(t, n)}。 這就是這題的
\textbf{暴力解}（brute force）：照題目的字面意思做，不動任何腦筋，所以幾乎不會錯。

這個程式要做的事，跟答案一樣多：答案是多少，\texttt{cnt++} 就執行幾次。 片段 1
在 $n = 10^9$ 時的答案大約是 $2 \times 10^9$，在本書測試的環境下這支程式跑一筆要
$1.1$ 秒左右。 $n$ 到 $10^{18}$ 時是它的十億倍，而且最多有 $10^5$ 筆，完全不可能。
所以要找一個不必真的跑完、就能算出次數的方法。

\textbf{片段 2 和片段 3：只跑外層。} 片段 2 的外層 \texttt{i} 每次乘 $2$，片段 3 的
外層每次除以 $3$，外層最多跑六十幾輪。 內層的次數直接寫得出來：片段 2 是 \texttt{i}
次，片段 3 是 \texttt{i + 1} 次（\texttt{j} 從 $0$ 開始）。 所以只要跑外層，每一輪
把內層的次數加上去。

\textbf{片段 1：換一個方向數。} 片段 1 的外層是 \texttt{i++}，要跑 $n$ 輪，只跑外層
也來不及。 先看 $n = 6$ 時每一輪的 \texttt{j}：

\begin{center}
\begin{tikzpicture}[x=1cm, y=1cm]
\foreach \i in {1,...,6} {\node[idxlbl] at (\i,0.75) {\texttt{i} = \i};}
\node[left, font=\small] at (0.4,0) {乘 $0$ 次};
\node[left, font=\small] at (0.4,-1) {乘 $1$ 次};
\node[left, font=\small] at (0.4,-2) {乘 $2$ 次};
\foreach \i in {1,...,6} {\node[hlcell] at (\i,0) {\i};}
\foreach \i/\j in {1/2,2/4,3/6} {\node[hlcell] at (\i,-1) {\j};}
\foreach \i in {4,5,6} {\node[cell, draw=black!20] at (\i,-1) {};}
\node[hlcell] at (1,-2) {4};
\foreach \i in {2,...,6} {\node[cell, draw=black!20] at (\i,-2) {};}
\node[font=\small] at (7.4,0) {$6$ 格};
\node[font=\small] at (7.4,-1) {$3$ 格};
\node[font=\small] at (7.4,-2) {$1$ 格};
\end{tikzpicture}
\diagcaption{片段 1 在 $n = 6$ 時，每一格是 \texttt{cnt++} 執行時 \texttt{j} 的值。
直著看是一輪外層，橫著看是「乘了幾次 $2$」。}
\label{fig:piece1-grid}
\end{center}

題目問的是格子總數。 直著一行一行數（照程式的順序）是 $3 + 2 + 2 + 1 + 1 + 1 = 10$，
要數 $n$ 行。 改成橫著一列一列數：第 $k$ 列是 \texttt{j} 已經乘了 $k$ 次 $2$ 的那
些格子，\texttt{j} $= i \times 2^k$。 這一格存在，剛好是 $i \times 2^k \le n$，也就
是 $i \le n / 2^k$，所以第 $k$ 列有 $\lfloor n / 2^k \rfloor$ 格。 總數是
\[
n + \left\lfloor \frac{n}{2} \right\rfloor + \left\lfloor \frac{n}{4} \right\rfloor + \cdots,
\]
加到 $\lfloor n / 2^k \rfloor = 0$ 為止。 $n = 6$ 時是 $6 + 3 + 1 = 10$，和直著數一
樣。 這個和最多六十幾項，用一個 \texttt{m /= 2} 的迴圈就能算出來。 而且
$\lfloor n / 2^k \rfloor$ 也等於 $\lfloor \lfloor n / 2^{k-1} \rfloor / 2 \rfloor$，
所以 \texttt{m} 從 $n$ 開始一直除以 $2$ 就好，不用算 $2^k$。

順便看出一件事：這個和小於 $n + n/2 + n/4 + \cdots = 2n$。 片段 1 雖然有兩層迴圈，
\texttt{cnt++} 只執行不到 $2n$ 次，跟 $n^2$ 差得很遠。

\textbf{程式}

\begin{example}[label=ex:loop-count]{cpp}{(loop\_count.cpp) 迴圈到底跑幾次}
#include <iostream>
using namespace std;

// 片段 1：第 k 列的 j 是 i 乘上 2 的 k 次方，
// 這一列有 n / 2^k 格（整數除法）
long long piece1(long long n) {
    long long total = 0;
    for (long long m = n; m > 0; m /= 2) {
        total += m;
    }
    return total;
}

// 片段 2：外層 i = 1, 2, 4, ...，內層各跑 i 次
long long piece2(long long n) {
    long long total = 0;
    for (long long i = 1; i <= n; i *= 2) {
        total += i;
    }
    return total;
}

// 片段 3：外層 i = n, n/3, n/9, ...，內層 j 從 0 到 i，跑 i + 1 次
long long piece3(long long n) {
    long long total = 0;
    for (long long i = n; i > 0; i /= 3) {
        total += i + 1;
    }
    return total;
}

int main() {
    ios::sync_with_stdio(false);
    cin.tie(nullptr);
    int q;
    cin >> q;
    while (q--) {
        int t;
        long long n;
        cin >> t >> n;
        long long ans;
        if (t == 1) {
            ans = piece1(n);
        } else if (t == 2) {
            ans = piece2(n);
        } else {
            ans = piece3(n);
        }
        cout << ans << '\n';
    }
    return 0;
}
\end{example}

\iosample{
5 \\
3 1 \\
2 7 \\
2 8 \\
3 9 \\
1 10 \\
}{
2 \\
7 \\
15 \\
16 \\
18 \\
}

\begin{itemize}
\item 把同一組輸入拿給暴力解跑，輸出也相同。 兩支程式做法完全不同，答案卻一樣，這本
身就是一種檢查；第~\ref{sec:testing}~節會把這個想法變成固定的流程。
\item 三個函式的迴圈都和原本的片段很像，差別只在內層被換成一個加法。 這樣寫比較容
易和原本的片段逐行對照。
\item 答案最大大約是 $2 \times 10^{18}$，在 \texttt{long long} 的範圍內（上限大約
$9.2 \times 10^{18}$）；\texttt{int} 只到大約 $2.1 \times 10^9$，一定不夠。
\end{itemize}

\textbf{複雜度}

每個問題的三個函式各最多跑 $\log_2 10^{18} \approx 60$ 輪，$10^5$ 個問題總共幾百
萬次加法。 最大的測試資料在本書測試的環境下約 $0.01$ 秒。

\subsection{常見錯誤}

\begin{itemize}
\item \textbf{看到兩層就套公式}：把片段 1 當成「外層 $n$ 次、內層 $\log_2 n$ 次」，寫成
$n \times (\lfloor \log_2 n \rfloor + 1)$。 範例 \texttt{1 10} 會印出 \texttt{40}，正確
是 \texttt{18}。 內層的 \texttt{j} 從 \texttt{i} 開始而不是從 $1$ 開始，\texttt{i}
越大，內層越短。
\item \textbf{漏掉從 $0$ 開始的那一次}：片段 3 寫成 \texttt{total += i}，範例
\texttt{3 1} 會印出 \texttt{1}，正確是 \texttt{2}。 \texttt{j} 從 $0$ 跑到
\texttt{i} 是 $i + 1$ 個數。
\item \textbf{用浮點數算次數}：片段 2 的外層跑 $\lfloor \log_2 n \rfloor + 1$ 輪，有人
用 \texttt{log2((double)n)} 算出輪數，再套 $2^{m+1} - 1$。 輸入 \texttt{2
576460752303423487}（$n = 2^{59} - 1$）時，正確答案是
\texttt{576460752303423487}，這樣寫在本書測試的環境下印出
\texttt{1152921504606846975}，多算了一輪。 原因在第~\ref{sec:overflow}~節：
\texttt{double} 存不下這麼大的整數，$2^{59} - 1$ 被捨入成 $2^{59}$。
\item \textbf{直接跑片段 1}：只把片段 2、3 改成只跑外層，片段 1 還是照原樣跑。 範例都
會對，但 $n = 10^{18}$ 的資料永遠跑不完。
\end{itemize}

\section{從次數到成長速度}
\label{sec:growth}

回到章首的診所。 第一種方法翻單子 $n(n - 1)/2$ 次。 第二種方法每張單子插一次、拿
一次，插一張算翻十張的時間，總共相當於翻 $10n + n = 11n$ 張。 兩種方法在幾個人數下
的比較：

\begin{center}
\begin{tabular}{rrrl}
\toprule
人數 $n$ & 第一種 $n(n-1)/2$ & 第二種 $11n$ & 比較省事的 \\
\midrule
$5$ & $10$ & $55$ & 第一種 \\
$20$ & $190$ & $220$ & 第一種 \\
$23$ & $253$ & $253$ & 一樣 \\
$24$ & $276$ & $264$ & 第二種 \\
$100$ & $4950$ & $1100$ & 第二種 \\
$1000$ & $499500$ & $11000$ & 第二種 \\
$10^5$ & 約 $5 \times 10^9$ & $1.1 \times 10^6$ & 第二種 \\
\bottomrule
\end{tabular}
\end{center}

人少的時候，第一種方法比較省事，因為它不用花插單子的那段固定成本。 兩邊在 $23$
人時剛好打平，從 $24$ 人起第二種方法一直比較快，而且差距越拉越大：一千人時差了
四十多倍，十萬人時差了四千多倍。

\begin{center}
\begin{tikzpicture}[x=0.25cm, y=0.0075cm]
\draw[->, black!60] (0,0) -- (42,0) node[right, font=\small] {$n$};
\draw[->, black!60] (0,0) -- (0,800) node[above, font=\small] {翻單子的次數};
\foreach \t in {10,20,30,40} {\draw[black!60] (\t,0) -- (\t,-15) node[below, font=\scriptsize] {\t};}
\foreach \t in {200,400,600} {\draw[black!60] (0,\t) -- (-0.4,\t) node[left, font=\scriptsize] {\t};}
\draw[thick, optStroke, domain=0:40, samples=60] plot (\x, {\x*(\x-1)/2});
\draw[thick, coreStroke, domain=0:40] plot (\x, {11*\x});
\draw[dashed, black!50] (23,0) -- (23,253);
\fill[black] (23,253) circle (2.5pt);
\node[font=\small, optStroke] at (29,620) {第一種};
\node[font=\small, coreStroke] at (39,350) {第二種};
\node[font=\scriptsize, above left] at (23,253) {$n = 23$};
\end{tikzpicture}
\diagcaption{兩種叫號方法翻單子的次數。 交叉點左邊第一種較少，右邊第二種較少。}
\label{fig:crossover}
\end{center}

如果插單子更慢，例如一張算翻一百張，交叉點會往右移到兩百人左右，但右邊最後還是第
二種贏。 理由是兩個次數長大的方式不一樣：第二種和 $n$ 成正比，人數變十倍，次數也
變十倍；第一種大約是 $n^2/2$，人數變十倍，次數變一百倍。 只要 $n$ 夠大，前面的係數
$10$ 或 $100$ 都追不上這個差距。

\subsection{只看最大的那一項}

把上面的觀察寫成一般的規則：$n$ 很大時，決定次數的是長得最快的那一項，係數和比較
小的項都不重要。 $n(n - 1)/2 = n^2/2 - n/2$，最大的項是 $n^2/2$，去掉係數就是
$n^2$，記作 $O(n^2)$；$11n$ 記作 $O(n)$。 這種寫法叫\textbf{時間複雜度}（time
complexity），讀作「big-O」。 白話的意思是：$n$ 夠大以後，次數不超過括號裡那個式
子的某個固定倍數。

\begin{itemize}
\item 兩段程式先後執行，複雜度取比較大的那段。 先排序 $O(n \log n)$、再掃一遍
$O(n)$，合起來是 $O(n \log n)$。
\item 一段 $O(f)$ 的程式被執行 $g$ 次，合起來是 $O(f \times g)$。
\item 次數會隨輸入的內容變（例如找到就 \texttt{break}）時，看最壞的情況。
\item $\log$ 的底數只差一個固定倍數（$\log_3 n$ 大約是 $\log_2 n$ 的 $0.63$ 倍），所以
$O(\log n)$ 不寫底數。
\end{itemize}

要分清楚兩種問題。 「這段程式的 \texttt{cnt++} 執行幾次」要的是精確的次數，例如
$n(n - 1)/2$，上一節的例題和本章練習的追蹤題都是這種。 「這個做法來不來得及」只要
量級，例如 $O(n^2)$，接下來估時間用的都是這種。 兩者都要會：精確的次數算錯了，量
級常常也跟著錯。

\subsection{實際量一量}

量級只說次數怎麼長大，沒說一秒能做幾次。 拿一個 $O(n^2)$ 的雙層迴圈來量：
對每一對 $i < j$ 檢查 \texttt{a[i] + a[j]} 是否等於某個數，共比較 $n(n - 1)/2$ 次。
在本書測試的環境下，三種大小各跑一次：

\begin{center}
\begin{tabular}{rrr}
\toprule
$n$ & 比較次數 $n(n-1)/2$ & 時間 \\
\midrule
$10^4$ & 約 $5 \times 10^7$ & $0.02$ 秒 \\
$3 \times 10^4$ & 約 $4.5 \times 10^8$ & $0.16$ 秒 \\
$10^5$ & 約 $5 \times 10^9$ & $1.8$ 秒 \\
\bottomrule
\end{tabular}
\end{center}

$n$ 變三倍，時間大約變九倍；$n$ 變十倍，時間大約變一百倍，和 $O(n^2)$ 的預測一致。
照這個比例，$n = 2 \times 10^5$ 會要七秒左右，大部分題目的時間限制是一兩秒，來不及。

一秒能做幾次，要看迴圈裡做的是什麼。 同一台機器上，另外量了兩種比較慢的操作：

\begin{center}
\begin{tabular}{lr}
\toprule
迴圈裡做的事 & 每秒大約幾次 \\
\midrule
加法和比較，依序讀一個陣列（上面的迴圈） & $3 \times 10^9$ \\
在一個 64 MB 的陣列裡，照亂數跳著讀 & $10^8$ \\
在約 60 萬個鍵的 \texttt{map} 裡插入或查詢 & $2 \times 10^6$ \\
\bottomrule
\end{tabular}
\end{center}

最快和最慢差了上千倍。 依序讀陣列時，電腦可以一次搬一大塊資料、同時做好幾個比較；
跳著讀時，每次都要等資料從記憶體送過來；\texttt{map} 的每次操作本身就要跳十幾二十
次。 銜接書 10.1 節的表用「一秒大約 $10^8$ 次簡單運算」來估，這個數字對單純的迴圈
很保守，但加上這些慢的操作之後剛好合用。 本書也用它：算出來在 $10^8$ 以下，一定
來得及；到 $10^9$ 附近，要看迴圈裡做的事有多單純；超過 $10^{10}$，換做法。

\subsection{從資料範圍反推需要的複雜度}
\label{sec:from-limits}

題目的限制欄寫著 $n \le 2 \times 10^5$，這個數字除了告訴你陣列要開多大，也在告訴你
需要多快的做法。 銜接書 10.1 節有一張「$n$ 的上限 $\to$ 可以接受的複雜度」的表，
本書不重印，只把用法整理成步驟：

\begin{enumerate}
\item 找出所有會讓工作量變大的數：資料筆數 $n$、詢問數 $q$、測試資料的組數、數值的
範圍。
\item 寫出你的做法的複雜度，把\textbf{每一個}數的上限都代進去。 詢問 $q$ 次、每次
$O(n)$，就是 $n \times q$，不是 $n$。
\item 和 $10^8$ 比。
\end{enumerate}

第二步最容易漏。 下面是幾個例子：

\begin{center}
\begin{tabular}{p{0.32\linewidth}p{0.22\linewidth}p{0.36\linewidth}}
\toprule
限制 & 想用的做法 & 估計 \\
\midrule
$n \le 2 \times 10^5$ & 排序後掃一遍 & $n \log_2 n \approx 3.6 \times 10^6$，可以 \\
$n, q \le 2 \times 10^5$ & 每個詢問掃一遍陣列 & $n \times q = 4 \times 10^{10}$，不行 \\
$n \le 5000$ & 檢查每一對 & $n^2/2 \approx 1.3 \times 10^7$，可以 \\
組數 $\le 10^5$，所有組的 $n$ 總和 $\le 10^6$ & 每組 $O(n)$ & 加起來 $O(10^6)$，可以 \\
同上，但每組都清空一個大小 $10^6$ 的陣列 & & $10^5 \times 10^6 = 10^{11}$，不行 \\
$n \le 20$ & 枚舉所有子集合 & $2^{20} \approx 10^6$，可以 \\
\bottomrule
\end{tabular}
\end{center}

倒數第二列是第~\ref{sec:ex-vote}~小節例題的重點：題目給「總和」的限制時，每一組
的工作量也要跟那一組的大小成正比，不能跟陣列宣告的大小成正比。

反過來，限制也會透露出題者想要的做法。 $n \le 20$ 幾乎是在說「可以試遍所有子集
合」；$n \le 5000$ 通常是 $O(n^2)$；$n$ 到 $10^5$ 以上就要 $O(n \log n)$ 或 $O(n)$；
$n$ 到 $10^{18}$ 時連 $O(n)$ 都不行，像上一節的例題，只能找公式或 $O(\log n)$ 的做法。

\begin{apcsnote}
實作題的題目會寫出資料範圍。 如果題目把測試資料分成幾個子題，範圍小的子題通常是給
比較直接的做法。 想不到夠快的做法時，先寫一個暴力解把小範圍的子題拿下來；這個暴力
解等一下也能拿來檢查你的快速版本（第~\ref{sec:testing}~節）。
\end{apcsnote}

\section{整數溢位與浮點數}
\label{sec:overflow}

\texttt{int} 能存的範圍大約是正負 $2.1 \times 10^9$，\texttt{long long} 大約是正負
$9.2 \times 10^{18}$。 有號整數算出超過範圍的值，叫\textbf{溢位}（overflow）。 在
C++ 裡這是未定義行為：程式不會停下來，也不會有錯誤訊息，只是印出一個錯的數字，而
且同一支程式換個編譯選項，錯的數字可能還不一樣。 銜接書 10.1 節的常見錯誤示範過
總和用 \texttt{int} 存的情形，這一節看幾種比較隱蔽的。

\subsection{先估最大值}

動手之前先算答案和中間值最大可能多大。 常見的幾種：

\begin{center}
\begin{tabular}{lll}
\toprule
算的東西 & 最大大約是 & 例子（$n = 2 \times 10^5$，值 $\le 10^9$） \\
\midrule
$n$ 個數的總和 & $n \times$ 最大值 & $2 \times 10^{14}$ \\
兩個數的乘積 & 最大值的平方 & $10^{18}$ \\
有幾對 & $n^2/2$ & $2 \times 10^{10}$ \\
$1 + 2 + \cdots + n$ & $n^2/2$ & $2 \times 10^{10}$ \\
\bottomrule
\end{tabular}
\end{center}

超過 $2 \times 10^9$ 就用 \texttt{long long}。 超過 $9 \times 10^{18}$ 時連
\texttt{long long} 都不夠，要換個算法，例如下面的「用除法比大小」。

\subsection{型別的上下限與型別別名}

範圍不用背。 \texttt{<limits>} 裡的 \texttt{numeric\_limits<T>::max()} 和
\texttt{min()} 給出型別 \texttt{T} 能存的最大值和最小值。 另外，\texttt{long long}
要一直寫很長，解題程式常在開頭寫 \texttt{using ll = long long;}，之後 \texttt{ll}
就是 \texttt{long long} 的另一個名字。 這叫\textbf{型別別名}（type alias），和 C 的
\texttt{typedef long long ll;} 意思完全一樣；\texttt{using} 的寫法是「新名字 =
舊型別」，從左讀到右，型別比較長時比較好讀。

\begin{example}[label=ex:limits]{cpp}{(limits.cpp) 型別的上下限與「無限大」}
#include <iostream>
#include <limits>
using namespace std;

using ll = long long;                   // 之後寫 ll 就是 long long

int main() {
    cout << numeric_limits<int>::max() << '\n';
    cout << numeric_limits<int>::min() << '\n';
    cout << numeric_limits<ll>::max() << '\n';
    const ll INF = 1e18;                // 當成「無限大」用
    cout << INF + INF << '\n';          // 兩個相加還在範圍內
    return 0;
}
\end{example}

\iosample{}{
2147483647 \\
-2147483648 \\
9223372036854775807 \\
2000000000000000000 \\
}

求最小值時，常先把答案設成「無限大」，再一個一個比。 直覺上會想用型別的最大值，可是之後只要拿它加任何正數就溢位了。 例如最短
路徑常寫「目前距離加上一條邊的長度」：

\begin{code}{cpp}{拿型別最大值當無限大（錯誤示範）}
ll dist = numeric_limits<ll>::max();   // 還沒走到
ll cand = dist + w;                    // w = 5
if (cand < dist) { ... }               // 以為找到更短的路
\end{code}

在本書測試的環境下，\texttt{cand} 變成 \texttt{-9223372036854775804}，比
\texttt{dist} 小，程式以為找到了一條更短的路。 所以「無限大」通常改用
\texttt{1e18} 這種常數：它比任何真正的答案都大，兩個相加也還在 \texttt{long long} 的
範圍內。 \texttt{1e18} 本身是 \texttt{double}，但 $10^{18}$ 能被 \texttt{double} 精確
表示，存進 \texttt{ll} 不會差。 \texttt{int} 的版本常用 \texttt{1e9}。

\subsection{結果存進 \texttt{long long} 也沒用}

下面這段程式的變數都宣告成 \texttt{long long} 了，可是三行裡有兩行是錯的：

\begin{code}[label=code:overflow-demo]{cpp}{三種在算式裡發生的溢位}
int w, h, n, b;
cin >> w >> h >> n >> b;               // 輸入 100000 100000 80000 36
long long area = w * h;                // 錯
long long area2 = (long long)w * h;    // 對
long long pairs = n * (n - 1) / 2;     // 錯
long long pairs2 = (long long)n * (n - 1) / 2;
long long mask = (1 << b) - 1;         // 錯
long long mask2 = (1LL << b) - 1;      // 對
\end{code}

在本書測試的環境下，用 \texttt{-O2} 編譯，印出這六個變數是：

\begin{code}{text}{上面六個變數的值}
area  = 1410065408     area2  = 10000000000
pairs = 1052476352     pairs2 = 3199960000
mask  = 15             mask2  = 68719476735
\end{code}

C++ 算 \texttt{w * h} 時只看兩個運算元的型別：兩個都是 \texttt{int}，就用
\texttt{int} 乘，乘完已經溢位，才轉成 \texttt{long long} 存起來。 等號左邊是什麼型
別，影響不到右邊怎麼算。 修正的方法是在相乘之前，先把其中一個運算元轉成
\texttt{long long}，另一個會自動跟著轉。 下面兩種寫法要分清楚：

\begin{itemize}
\item \texttt{(long long)w * h}：先轉 \texttt{w}，再乘，對。
\item \texttt{(long long)(w * h)}：先用 \texttt{int} 乘完，再轉，跟沒轉一樣。
\end{itemize}

\texttt{1 << b} 也是同一回事：\texttt{1} 是 \texttt{int}，左移 $36$ 位超過了
\texttt{int} 的 $32$ 位。 要寫 \texttt{1LL << b}，\texttt{LL} 讓這個 \texttt{1} 一開
始就是 \texttt{long long}。 這一行在加上 \texttt{-fsanitize=undefined} 編譯時印出
的是 \texttt{-1}，不是 \texttt{15}，這就是「未定義行為」的意思：結果沒有任何保證。

\subsection{用除法比大小}

有時候乘出來的值連 \texttt{long long} 都存不下，但我們只想知道它有沒有超過某個數。
例如想檢查 $n^2 \le T$，$T$ 到 $10^{18}$，$n$ 也可能到 $10^{18}$。 這時把乘法換成除
法：
\[
n \times n \le T \quad\text{等於}\quad n \le \left\lfloor \frac{T}{n} \right\rfloor
\qquad (n \ge 1).
\]
因為 $n$ 是整數，「$n \le T/n$」和「$n$ 不超過 $T/n$ 的整數部分」是同一件事。 右邊
是整數除法，結果不會比 $T$ 大，一定不會溢位。 三個數相乘也一樣，$n^3 \le T$ 可以
寫成 \texttt{n <= T / n / n}。

另一個常見的地方是兩個數的中點。 \texttt{lo} 和 \texttt{hi} 都是 \texttt{int}、可能
到 $2 \times 10^9$ 附近時，\texttt{(lo + hi) / 2} 的加法就溢位了。 改寫成
\texttt{lo + (hi - lo) / 2}，算出來的值一樣，但中間值不會超過 \texttt{hi}。 第~3~章
的二分搜尋會一直用到這個寫法。

\subsection{浮點數存不下大整數}

\texttt{double} 有 $53$ 個二進位位數存數字，能精確表示的整數只到
$2^{53} \approx 9 \times 10^{15}$。 再大的整數會被捨入成最接近的可表示值，例如
$2^{53} + 1 = 9007199254740993$ 轉成 \texttt{double} 再轉回來，會變成
$9007199254740992$。

這會讓「用 \texttt{sqrt}、\texttt{log2} 算整數答案」的程式在大數上差一：

\begingroup\sloppy
\begin{itemize}
\item $T = 10^{18} - 1$，想求最大的 $n$ 使 $n^2 \le T$，答案是 $999999999$。
\texttt{(long long)sqrt((double)T)} 卻得到 $1000000000$，因為 $10^{18} - 1$ 轉成
\texttt{double} 時被捨入成 $10^{18}$。
\item $n = 2^{59} - 1$，$\lfloor \log_2 n \rfloor = 58$。
\texttt{(long long)log2((double)n)} 得到 $59$，理由一樣。
\end{itemize}
\endgroup

浮點數還是有用處：拿它估一個大概的值，再用整數修正。 下一小節的例題會示範：
先用 \texttt{sqrt} 估，再用上一小節的除法檢查，大了就減一、小了就加一。

\subsection{例題：預算內的最大規模}
\label{sec:ex-max-size}

\textbf{題意}（\akiref{largest-size-within-the-budget}{預算內的最大規模}）

\begin{itemize}
\item 有六個函式：$f_1(n) = n$；$f_2(n) = n \times b(n)$，$b(n)$ 是 $n$ 的二進位位數，
$f_2(0) = 0$；$f_3(n) = n^2$；$f_4(n) = n^3$；$f_5(n) = 2^n$；$f_6(n) = n!$。
\item 有 $q$ 個詢問，每個給代號 $t$ 和預算 $T$，求最大的非負整數 $n$ 使
$f_t(n) \le T$。
\item $q \le 10^5$，$1 \le T \le 10^{18}$。
\end{itemize}

$T = 10^8$ 時六個答案依序是 $10^8$、$4347826$、$10^4$、$464$、$26$、$11$。 這一列
就是銜接書那張表的由來：一秒做 $10^8$ 次的話，$O(n^2)$ 的做法能處理一萬筆，$O(2^n)$
的做法只能處理二十幾筆。 $f_2$ 是 $n \log n$ 的一種精確版本。

\textbf{想法}

暴力解是 $n$ 從 $0$ 開始一個一個試，直到 $f(n) > T$。 $f_5$、$f_6$ 長得很快，試幾十
個就停了；可是 $f_1$ 在 $T = 10^{18}$ 時要試 $10^{18}$ 次，$f_3$ 要試 $10^9$ 次，再乘上
$10^5$ 個詢問，不可能。 六個函式分開處理：

\begin{itemize}
\item \textbf{$f_1$}：答案就是 $T$。
\item \textbf{$f_5$、$f_6$}：照暴力解一個一個試，最多六十步。 判斷「下一個還在預算內」
時用除法：$2^{n+1} \le T$ 寫成 \texttt{value <= T / 2}，$(n+1)! \le T$ 寫成
\texttt{value <= T / (n + 1)}。 這兩題其實直接乘也不會溢位，但要確定這件事得自己算
一遍；用除法就不用算。
\item \textbf{$f_3$、$f_4$}：用 \texttt{sqrt}、\texttt{cbrt} 估一個 $n$，再用除法修正：
$n^2 > T$ 就往下減，$(n+1)^2 \le T$ 就往上加。 浮點數只差一兩個，修正的迴圈只跑幾
次。
\item \textbf{$f_2$}：依位數分組。 有 $b$ 位的數是 $2^{b-1}$ 到 $2^b - 1$，這一組裡
$f_2(n) = n \times b$，條件 $n \times b \le T$ 等於 $n \le \lfloor T / b \rfloor$，
所以這一組最大的答案是 $\min(2^b - 1, \lfloor T / b \rfloor)$，前提是它不小於
$2^{b-1}$（否則這一組沒有任何數在預算內）。 $f_2$ 隨 $n$ 遞增，所以位數越多、還能
在預算內的組，答案越大；$b$ 從 $1$ 試到 $62$，留下最後一個合格的就好。
\end{itemize}

$f_2$ 的上限為什麼試到 $62$？ $T \le 10^{18} < 2^{60}$，答案不超過 $T$，最多 $60$
位；$62$ 只是留一點餘裕，而且 \texttt{1LL << 62} 還在 \texttt{long long} 的範圍內。

\textbf{程式}

\begin{example}[label=ex:max-size]{cpp}{(max\_size.cpp) 預算內的最大規模}
#include <cmath>
#include <iostream>
using namespace std;

// 每個函式回答：最大的 n，使 f(n) <= T。比大小一律用除法，不做會溢位的乘法。
// f(n) = n 的答案就是 T，直接寫在 main 裡。

long long nLogN(long long T) {              // f(n) = n * (n 的二進位位數)
    long long best = 0;
    for (int b = 1; b <= 62; b++) {         // 有 b 位的 n：2^(b-1) 到 2^b - 1
        long long lo = 1LL << (b - 1), hi = (1LL << b) - 1;
        long long n = min(hi, T / b);       // n * b <= T  等於  n <= T / b
        if (n >= lo) {
            best = n;
        }
    }
    return best;
}

// n^k <= T 嗎？（n >= 1）把 T 連除 k 次 n，還剩至少 1 就是
bool fits(long long n, int k, long long T) {
    for (int i = 0; i < k; i++) {
        T /= n;
    }
    return T >= 1;
}

long long root(long long T, int k) {        // f(n) = n^k，k = 2 或 3
    double guess = (k == 2) ? sqrt((double)T) : cbrt((double)T);
    long long n = (long long)guess;         // 先用浮點數估一個大概
    while (n > 0 && !fits(n, k, T)) {       // 太大，往下修
        n--;
    }
    while (fits(n + 1, k, T)) {             // 還能再大，往上修
        n++;
    }
    return n;
}

long long power2(long long T) {             // f(n) = 2^n
    long long n = 0, value = 1;             // value = 2^n
    while (value <= T / 2) {                // 2^(n+1) <= T
        value *= 2;
        n++;
    }
    return n;
}

long long factorial(long long T) {          // f(n) = n!
    long long n = 1, value = 1;             // value = n!，0! = 1! = 1
    while (value <= T / (n + 1)) {          // (n+1)! <= T
        n++;
        value *= n;
    }
    return n;
}

int main() {
    ios::sync_with_stdio(false);
    cin.tie(nullptr);
    int q;
    cin >> q;
    while (q--) {
        int t;
        long long T;
        cin >> t >> T;
        long long ans;
        if (t == 1) {
            ans = T;
        } else if (t == 2) {
            ans = nLogN(T);
        } else if (t <= 4) {
            ans = root(T, t - 1);
        } else if (t == 5) {
            ans = power2(T);
        } else {
            ans = factorial(T);
        }
        cout << ans << '\n';
    }
    return 0;
}
\end{example}

\iosample{
6 \\
1 100000000 \\
2 100000000 \\
3 100000000 \\
4 100000000 \\
5 100000000 \\
6 100000000 \\
}{
100000000 \\
4347826 \\
10000 \\
464 \\
26 \\
11 \\
}

\begin{itemize}
\item \texttt{fits} 把 $n^k \le T$ 換成「$T$ 連除 $k$ 次 $n$ 之後至少是 $1$」，和上一
小節的除法是同一個想法。 \texttt{root} 的第一個 \texttt{while} 要先檢查
\texttt{n > 0}，因為除數不能是 $0$。
\item \texttt{factorial} 從 $n = 1$ 開始，因為 $0! = 1! = 1 \le T$，答案至少是 $1$。
\texttt{power2} 從 $n = 0$ 開始：$T = 1$ 時 $2^1 = 2 > 1$，答案是 $0$。
\item \texttt{nLogN} 裡的 \texttt{1LL << b} 不能寫成 \texttt{1 << b}，$b \ge 31$ 時就
是上一小節的位移溢位。
\end{itemize}

\textbf{複雜度}

每個詢問最多六十幾步（$f_2$ 的 $62$ 組、$f_5$ 的六十次倍增），修正迴圈只跑幾次。
$10^5$ 個詢問大約 $6 \times 10^6$ 步，最大的測試資料在本書測試的環境下約 $0.01$ 秒。

\subsection{常見錯誤}

\begingroup\sloppy
\begin{itemize}
\item \textbf{只用浮點數}：$f_3$ 寫成 \texttt{(long long)sqrt((double)T)} 就回傳，輸入
\texttt{3 999999999999999999} 會印出 \texttt{1000000000}，正確是 \texttt{999999999}。
$f_4$ 用 \texttt{cbrt} 也一樣，同一個 $T$ 印出 \texttt{1000000}，正確是
\texttt{999999}。 $f_5$ 用 \texttt{log2}，輸入 \texttt{5 576460752303423487} 會印出
\texttt{59}，正確是 \texttt{58}。 範例的 $T$ 都很小，這三種寫法都會通過範例。
\item \textbf{比大小時直接乘}：\texttt{fits} 寫成把 $n$ 連乘 $k$ 次再和 $T$ 比。 這裡
$n$ 只到 $10^9$ 左右，乘起來還在範圍內；但同樣的寫法搬到 $f_2$ 的 \texttt{hi * b} 上，
$b = 62$ 時 \texttt{hi} 大約 $4.6 \times 10^{18}$，乘上 $62$ 就溢位了。 與其每次都去確認會不會溢位，不如一律用
除法。
\item \textbf{位數算少一位}：把 $b(n)$ 當成 $\lfloor \log_2 n \rfloor$，也就是從 $0$
開始數位數。 範例的第二列會印出 \texttt{4545454}，比正確的 \texttt{4347826} 大。
$8 = 1000_2$ 有 $4$ 位，不是 $3$ 位。
\item \textbf{小的邊界}：忘了 $0! = 1$，或以為 $T = 1$ 時 $f_5$ 的答案是 $1$。 這些
在範例 2 裡都有，交出去前把範例全部跑過一次。
\end{itemize}
\endgroup

\section{輸入與輸出}
\label{sec:io}

演算法再快，資料讀得太慢一樣會超過時間限制；答案算對了，多印一個空白也可能被判錯。
這一節講三件事：讀寫的速度、讀到檔案結束、多組資料。

\subsection{讀寫的速度}

\begin{bridgerecap}[輸入加速的兩行]
\texttt{main} 一開頭加上
\texttt{ios::sync\_with\_stdio(false);} 和 \texttt{cin.tie(nullptr);}（銜接書第~1~章）。
第一行關掉 \texttt{cin}／\texttt{cout} 和 C 的 \texttt{scanf}／\texttt{printf} 之間的同
步；第二行讓 \texttt{cin} 讀資料前不再先把 \texttt{cout} 的緩衝區送出去。 加了之後
不要再混用 \texttt{scanf}／\texttt{printf}，否則輸出的順序可能會亂掉。 換行用
\texttt{'\textbackslash n'}，不要用 \texttt{endl}：\texttt{endl} 除了換行，還會把緩
衝區立刻送出去，每一行都送一次就很慢。
\end{bridgerecap}

這三件事的效果，在本書測試的環境下量了一次。 程式讀 $10^6$ 個整數，每讀一個就印一
行，輸入和輸出都是檔案：

\begin{center}
\begin{tabular}{lr}
\toprule
寫法 & 時間 \\
\midrule
什麼都不加，換行用 \texttt{'\textbackslash n'} & $0.70$ 秒 \\
只加 \texttt{sync\_with\_stdio(false)} & $0.44$ 秒 \\
兩行都加，但換行用 \texttt{endl} & $0.47$ 秒 \\
兩行都加，換行用 \texttt{'\textbackslash n'} & $0.10$ 秒 \\
\bottomrule
\end{tabular}
\end{center}

少做任何一件，都會慢四倍以上。 只讀不印的時候，讀 $2 \times 10^6$ 個整數，不加速要
$0.33$ 秒，加了兩行是 $0.06$ 秒，和 \texttt{scanf} 的 $0.07$ 秒差不多。 一題的時間
限制常常只有一秒，這幾倍的差距可能就是過與不過的差別。 所以本書每一支讀大量資料的程
式，開頭都有這兩行。

\subsection{讀到檔案結束}

有些題目不告訴你有幾筆資料，資料一直到檔案結束為止。 \texttt{cin >> x} 這個運算式
本身有值：讀到東西時當成 \texttt{true}，讀不到（檔案結束，或內容不是數字）時當成
\texttt{false}。 所以 \texttt{while (cin >> x)} 會一直讀，讀到檔案結束就停。

\begin{code}[label=code:delivery]{cpp}{讀到檔案結束：不知道有幾筆的外送距離}
long long d, total = 0;
int count = 0;
while (cin >> d) {              // 讀不到數字（檔案結束）時停下
    count++;
    total += d;
}
cout << count << ' ' << total << '\n';
\end{code}

每個數字是一筆外送的距離（公尺），輸入 \texttt{1200}、\texttt{3500 800}、
\texttt{2750}（三行）時印出 \texttt{4 8250}。 \texttt{>>} 會跳過所有空白和換行，所
以一行有幾個數字都沒關係，檔案最後有沒有換行也沒關係。 在鍵盤上手動輸入時，要按
Ctrl+D（Windows 是 Ctrl+Z 再按 Enter）表示輸入結束；用 \texttt{./a.out < in.txt}
從檔案讀時，讀完檔案自然就結束了。

如果資料是「一行一筆」，而一行裡的個數不固定，就要知道哪些數字在同一行。 這時用
\texttt{getline} 讀一整行，再用 \texttt{istringstream} 把這一行當成一個小的輸入來拆。
\texttt{istringstream} 在 \texttt{<sstream>} 裡，用法和 \texttt{cin} 一樣：

\begin{example}[label=ex:orders]{cpp}{(orders.cpp) 每位客人點了幾樣}
#include <iostream>
#include <sstream>
#include <string>
using namespace std;

int main() {
    int n;
    cin >> n;
    string line;
    getline(cin, line);             // 讀掉 n 後面剩下的換行
    for (int i = 1; i <= n; i++) {
        getline(cin, line);         // 一整行：這位客人點的所有品項
        istringstream in(line);
        string item;
        int count = 0;
        while (in >> item) {        // 從這一行裡一個一個拆出品項
            count++;
        }
        cout << "#" << i << ' ' << count << '\n';
    }
    return 0;
}
\end{example}

\iosample{
3 \\
latte bagel \\
tea \\
mocha muffin scone \\
}{
\#1 2 \\
\#2 1 \\
\#3 3 \\
}

第 10 行的 \texttt{getline} 是必要的。 \texttt{cin >> n} 讀完 \texttt{3} 就停下來，後
面的換行還留在輸入裡；少了這一行，第一次 \texttt{getline} 會讀到那個空的行尾，程式
印出 \texttt{\#1 0}、\texttt{\#2 2}、\texttt{\#3 1}，整個錯開一行。 這是銜接書第~3~章
講過的陷阱，在混用 \texttt{>>} 和 \texttt{getline} 的題目裡一定會遇到。

\subsection{控制小數位數與寬度}

題目常要求「四捨五入到小數點後第二位」或「時間寫成 \texttt{09:05}」。 這些用
\texttt{<iomanip>} 裡的格式設定：\texttt{fixed} 加 \texttt{setprecision(k)} 印出小數
點後 $k$ 位；\texttt{setw(w)} 讓下一個輸出至少佔 $w$ 格，\texttt{setfill(c)} 決定空出
來的格子填什麼字元。

\begin{example}[label=ex:price]{cpp}{(price.cpp) 平均金額與補零的時間}
#include <iomanip>
#include <iostream>
using namespace std;

int main() {
    int n;
    cin >> n;
    long long total = 0;
    for (int i = 0; i < n; i++) {
        int x;
        cin >> x;
        total += x;
    }
    double avg = (double)total / n;
    cout << fixed << setprecision(2) << avg << '\n';   // 小數點後兩位
    int h, m;
    cin >> h >> m;
    cout << setfill('0') << setw(2) << h << ':' << setw(2) << m << '\n';
    return 0;
}
\end{example}

\iosample{
3 \\
1000 2050 1235 \\
9 5 \\
}{
1428.33 \\
09:05 \\
}

\begin{itemize}
\item \textbf{少了 \texttt{fixed}}：\texttt{setprecision(2)} 的意思變成「全部只留兩位
有效數字」，同一個平均值會印成 \texttt{1.4e+03}。
\item \textbf{設定會一直留著}：\texttt{fixed}、\texttt{setprecision}、\texttt{setfill}
設了之後對後面所有輸出都有效，之後再印 \texttt{3.0} 會變成 \texttt{3.00}。 只有
\texttt{setw} 是只管下一個輸出，所以上面的程式寫了兩次。
\item \textbf{四捨五入不一定照十進位}：\texttt{2.675} 在 \texttt{double} 裡其實比
$2.675$ 小一點，印到兩位是 \texttt{2.67}；剛好在中間的 \texttt{0.125} 也印成
\texttt{0.12}。 題目要求精確的四捨五入時，金額先用整數算（以分為單位），最後再轉。
\end{itemize}

同一套設定也能用在 \texttt{ostringstream} 上：把它當成 \texttt{cout} 來輸出，最後用
\texttt{.str()} 取出字串，就能先把格式化好的結果存起來，例如拼出一行再一次印出。

\subsection{例題：校慶攤位主題票選}
\label{sec:ex-vote}

\textbf{題意}（\akiref{voting-on-school-fair-booth-themes}{校慶攤位主題票選}）

\begin{itemize}
\item 輸入有很多組，每組是一個班級：第一行是選票張數 $n$ 和主題數 $k$；$n \ge 1$ 時下
一行有 $n$ 個主題編號（$1$ 到 $k$），$n = 0$ 時沒有這一行。
\item 最後一行是 \texttt{0 0}，表示輸入結束。
\item 每組輸出得票最多的主題編號，同票時取編號最小的；沒有人投票時輸出 $0$。
\item 組數最多 $5 \times 10^5$，每組的 $n$、$k$ 最多 $5 \times 10^5$，所有組的 $n + k$
加起來不超過 $10^6$。
\end{itemize}

例如輸入是 \texttt{3 5}、\texttt{5 5 2}、\texttt{0 4}、\texttt{1 2}、\texttt{2}、
\texttt{0 0}（每項一行）。 第一組主題 $5$ 得兩票，輸出 $5$；第二組有四個主題、沒有
人投票，輸出 $0$；第三組輸出 $2$。

\textbf{想法}

每一組的計票很直接：開一個計數陣列，\texttt{cnt[c]} 是主題 $c$ 的票數，讀一張票就加
一，最後從 $1$ 到 $k$ 找最大的。 這題要想清楚的是另外三件事。

\textbf{結束的條件。} 第二組的 \texttt{0 4} 是一組合法的資料：有四個主題，只是沒有人
投票。 只有 $n$ 和 $k$ 都是 $0$ 才是結束標記。 讀到 $n = 0$ 時也不用特別處理「沒有
選票那一行」，因為 \texttt{>>} 本來就會跳過換行，直接讀下一組的 $n$、$k$。

\textbf{每一組的工作量。} 用第~\ref{sec:from-limits}~小節的方法估。 如果宣告一個大小
$5 \times 10^5 + 1$ 的陣列，每組開始前整個清成 $0$，最壞的情況是 $5 \times 10^5$ 組
都要清一次，共 $2.5 \times 10^{11}$ 次，來不及。 題目限制的是所有組的 $n + k$ 總和，
所以每一組的工作量要和這一組的 $n + k$ 成正比。 最簡單的做法，是每一組各自開一個大小
$k + 1$ 的 \texttt{vector}：開的時候順便清成 $0$，花 $O(k)$；讀票 $O(n)$；找最大值
$O(k)$。 所有組加起來是 $O(\sum (n + k)) = O(10^6)$。

\textbf{同票和沒有票。} 從 \texttt{best = 0} 開始，主題 $c$ 的票數嚴格大於
\texttt{cnt[best]} 時才換成 $c$。 \texttt{cnt[0]} 永遠是 $0$，所以沒有人投票時
\texttt{best} 一直是 $0$，剛好是題目要的輸出；有人投票時，第一個票數大於 $0$ 的主題
會換上來。 用嚴格大於，同票時就會留下編號小的那個。

同樣的模型也出現在別的地方：客服系統每天統計哪一類問題被反映得最多，外送平台統計每
家店當天最多人點的品項。 資料一天一組、類別數每天不同，就是這題的形狀。

\textbf{程式}

\begin{example}[label=ex:vote]{cpp}{(vote.cpp) 校慶攤位主題票選}
#include <iostream>
#include <vector>
using namespace std;

int main() {
    ios::sync_with_stdio(false);
    cin.tie(nullptr);
    int n, k;
    while (cin >> n >> k) {
        if (n == 0 && k == 0) {           // 兩個都是 0 才是結束標記
            break;
        }
        vector<int> cnt(k + 1, 0);        // 這一組自己的計數，大小跟著 k
        for (int i = 0; i < n; i++) {
            int a;
            cin >> a;
            cnt[a]++;
        }
        int best = 0;                     // cnt[0] 一定是 0
        for (int c = 1; c <= k; c++) {
            if (cnt[c] > cnt[best]) {     // 嚴格大於才換，同票留小的編號
                best = c;
            }
        }
        cout << best << '\n';
    }
    return 0;
}
\end{example}

\iosample{
3 5 \\
5 5 2 \\
0 4 \\
1 2 \\
2 \\
0 0 \\
}{
5 \\
0 \\
2 \\
}

\begin{itemize}
\item \texttt{while (cin >> n >> k)} 同時處理兩種結束：讀到 \texttt{0 0}，或是檔案在那
之前就結束了。 題目保證有 \texttt{0 0}，多這一層保護不花任何成本。
\item \texttt{vector<int> cnt(k + 1, 0)} 每一組都重新建立，上一組的票數不會留下來，也
不用自己清。
\end{itemize}

\textbf{複雜度}

所有組加起來 $O(\sum (n + k))$，最多大約 $10^6$ 次操作，加上讀 $10^6$ 個整數。 最大
的測試資料在本書測試的環境下約 $0.03$ 秒。

\subsection{常見錯誤}

\begin{itemize}
\item \textbf{看到 $n = 0$ 就結束}：結束條件寫成 \texttt{if (n == 0) break;}，上面的範
例只會印出 \texttt{5}，第二組和第三組都漏掉了。
\item \textbf{每組清空整個大陣列}：宣告一個大小 $5 \times 10^5 + 1$ 的計數陣列，每組都
整個清成 $0$。 答案是對的，但組數很多的測試資料在本書測試的環境下要跑 $9$ 秒左右，
正確的程式只要 $0.03$ 秒。
\item \textbf{忘了清}：計數陣列宣告在迴圈外面，每一組都沒有清。 上面的範例會印出
\texttt{5}、\texttt{2}、\texttt{2}：第二組沒有人投票，卻讀到第一組留下的票數。
\item \textbf{同票用 \texttt{>=}}：比較寫成 \texttt{cnt[c] >= cnt[best]}，同票時會換成
編號大的。 輸入 \texttt{4 3}、\texttt{1 3 3 1}、\texttt{0 0} 會印出 \texttt{3}，正確
是 \texttt{1}。
\end{itemize}

\subsection{輸出的格式}

評分系統拿你的輸出和正確答案比對。 有些系統會忽略行尾的空白和最後多出來的空行，有
些逐字比對；照最嚴格的那種準備就不用猜。 數字之間的空白、換行、\texttt{Case \#1: 5}
這種格式的大小寫和冒號，都照題目寫；一行印好幾個數字時，把空白印在數字「前面」（第
一個除外），最後就不會多一個空白；題目沒要求的提示文字一個字都不要印；除錯用的輸出改
印到 \texttt{cerr}，它不會被比對，但交出去前還是刪掉，印太多也花時間。

\section{自己測試：範例之外}
\label{sec:testing}

題目給的範例是用來說明題意的，通常很小，也不會故意去碰每一個邊界。 本章前面的常見
錯誤裡，有好幾個錯法都能通過全部範例：用 \texttt{sqrt} 求整數平方根、片段 1 直接跑、
清空整個大陣列。 交出去之前，至少要做兩件事：自己補邊界的輸入，以及在大量的隨機小輸
入上，拿暴力解來比對。

\subsection{手造邊界輸入}

想邊界輸入時，從題目的每一個限制下手：

\begin{itemize}
\item \textbf{最小}：$n = 1$、$n = 0$（題目允許的話）、只有一組資料。
\item \textbf{最大}：每個數都取上限，檢查溢位和時間。 迴圈次數的例題要測
$n = 10^{18}$，主題票選要測「很多組、每組很小」和「一組、很大」兩種。
\item \textbf{剛好在分界上}：題目的條件用到 $\le$ 或 $<$ 的地方，取剛好相等的值。 例如
$T$ 剛好是某個完全平方數、比它少一，或 $n$ 剛好是 $2$ 的次方、比它少一。
\item \textbf{全部相同、已經排好、完全反過來}：很多錯誤只在這些情況出現。
\item \textbf{格式的邊界}：結束標記前一組是空的（\texttt{0 4}）、檔案最後沒有換行。
\end{itemize}

每一筆都要先手算出正確答案，再和程式比。 手算不出來的大輸入，至少看程式有沒有在時
間內跑完，答案的大小合不合理。

\subsection{用暴力解比對}

手造的輸入受限於你想得到的情況。 想不到的那些，交給電腦：寫一個慢但簡單的暴力解，
產生幾千組隨機的小輸入，看兩支程式的輸出是否完全相同。 這種做法俗稱
\textbf{對拍}。 需要三個程式和一個腳本：

\begin{center}
\begin{tikzpicture}[node distance=0.8cm and 0.8cm, font=\small]
\node[process] (gen) {產生器};
\node[ioNode, right=of gen] (in) {\texttt{test.in}};
\node[process, above right=0.2cm and 0.8cm of in] (mine) {要測的程式};
\node[process, below right=0.2cm and 0.8cm of in] (brute) {暴力解};
\node[decision, right=3.4cm of in, aspect=1.8, inner sep=1pt] (cmp) {相同？};
\draw[flowarrow] (gen) -- (in);
\draw[flowarrow] (in) -- (mine.west);
\draw[flowarrow] (in) -- (brute.west);
\draw[flowarrow] (mine.east) -- (cmp);
\draw[flowarrow] (brute.east) -- (cmp);
\draw[flowarrow] (cmp.south) -- ++(0,-1.3) -| node[pos=0.25, below] {相同：換下一個種子} (gen.south);
\node[right=0.4cm of cmp, align=left] (stop) {不同：印出\\\texttt{test.in}，停下};
\draw[flowarrow] (cmp) -- (stop);
\end{tikzpicture}
\diagcaption{用暴力解比對的流程。 每一輪用不同的亂數種子產生一組輸入。}
\label{fig:check-loop}
\end{center}

以主題票選為例。 暴力解要寫得和正式的程式越不一樣越好，這樣兩邊才不會犯同一個錯。
正式的程式用計數陣列，暴力解就對每個主題把整組票從頭數一次，$O(nk)$，很慢，但幾乎
不可能寫錯：

\begin{example}[label=ex:vote-brute]{cpp}{(vote\_brute.cpp) 主題票選的暴力解}
#include <iostream>
#include <vector>
using namespace std;

// 暴力解：每個主題都把整組票從頭數一次
int main() {
    int n, k;
    while (cin >> n >> k) {
        if (n == 0 && k == 0) {
            break;
        }
        vector<int> a(n);
        for (int &x : a) {
            cin >> x;
        }
        int best = 0, bestVotes = 0;
        for (int c = 1; c <= k; c++) {
            int votes = 0;
            for (int x : a) {
                if (x == c) {
                    votes++;
                }
            }
            if (votes > bestVotes) {
                best = c;
                bestVotes = votes;
            }
        }
        cout << best << '\n';
    }
    return 0;
}
\end{example}

\iosample{
5 3 \\
2 3 2 1 3 \\
4 4 \\
4 1 4 1 \\
0 0 \\
}{
2 \\
1 \\
}

產生器需要亂數。 C++ 用 \texttt{<random>} 裡的 \texttt{mt19937} 產生亂數：建立時給
一個\textbf{種子}（seed），之後每次呼叫 \texttt{rng()} 得到一個 $0$ 到 $2^{32} - 1$ 的
整數；再交給 \texttt{uniform\_int\_distribution<int>(lo, hi)}，就變成 \texttt{lo} 到
\texttt{hi} 之間平均分布的整數。 同一個種子一定得到同一串亂數：

\begin{example}[label=ex:rng]{cpp}{(rng\_demo.cpp) 固定種子的亂數}
#include <iostream>
#include <random>
using namespace std;

int main() {
    mt19937 rng(2026);                      // 種子固定，每次執行都一樣
    cout << rng() << ' ' << rng() << ' ' << rng() << '\n';  // 0 到 2^32 - 1
    uniform_int_distribution<int> die(1, 6);    // 1 到 6 平均分布
    for (int i = 0; i < 10; i++) {
        cout << (i > 0 ? " " : "") << die(rng);
    }
    cout << '\n';
    return 0;
}
\end{example}

\iosample{}{
942082305 3292560774 1773871898 \\
6 6 6 1 5 3 1 6 4 2 \\
}

第一行在任何編譯器上都一樣，因為 C++ 標準把 \texttt{mt19937} 的演算法規定死了；第二
行經過 \texttt{uniform\_int\_distribution}，標準沒有規定換算的細節，換一套標準函式庫
可能不同，但同一個環境裡一定重現得出來。 C 的 \texttt{rand()} 不適合用在這裡：它的
最大值 \texttt{RAND\_MAX} 在有些環境下只有 $32767$；\texttt{rand() \% k} 在 $k$ 不整除
範圍時，小的數字出現得比較多；種子還要另外用 \texttt{srand} 設，整支程式共用一個狀
態。

主題票選的產生器從命令列讀種子，產生一組合法的小輸入。 種子不同，輸入就不同；出錯時
用同一個種子就能重現：

\begin{example}[label=ex:vote-gen]{cpp}{(vote\_gen.cpp) 主題票選的隨機輸入產生器}
#include <iostream>
#include <random>
#include <string>
using namespace std;

// 用法：./vote_gen 種子
// 產生 1 到 3 組小資料，最後一行是 0 0
int main(int argc, char *argv[]) {
    if (argc < 2) {
        cerr << "usage: vote_gen seed\n";
        return 1;
    }
    mt19937 rng(stoul(argv[1]));
    auto randInt = [&](int lo, int hi) {   // lo 到 hi 之間的整數
        return uniform_int_distribution<int>(lo, hi)(rng);
    };
    int groups = randInt(1, 3);
    for (int g = 0; g < groups; g++) {
        int n = randInt(0, 5), k = randInt(1, 4);
        cout << n << ' ' << k << '\n';
        for (int i = 0; i < n; i++) {
            cout << randInt(1, k) << (i + 1 < n ? ' ' : '\n');
        }
    }
    cout << "0 0\n";
    return 0;
}
\end{example}

\begin{itemize}
\item \texttt{main} 的兩個參數是命令列引數：\texttt{argc} 是個數，\texttt{argv[1]} 是
程式名稱後面的第一個字串。 執行 \texttt{./vote\_gen 7} 時，\texttt{argv[1]} 是
\texttt{"7"}，用 \texttt{stoul} 轉成整數當種子。
\item 範圍故意設得很小：最多三組、每組最多五張票、四個主題。 小輸入出錯時容易手算，
而且同票、沒有票、\texttt{n = 0} 這些情況會經常出現。
\end{itemize}

最後是把它們串起來的腳本。 下面是 Linux 和 macOS 的 bash 寫法；Windows 可以在 WSL
或 Git Bash 裡執行同一個腳本：

\begin{code}[label=code:check-sh]{bash}{check.sh：用暴力解比對}
#!/bin/bash
# 用法：./check.sh 要測的程式 暴力解 產生器 次數
for ((i = 1; i <= $4; i++)); do
    ./$3 $i > test.in
    ./$1 < test.in > mine.out
    ./$2 < test.in > brute.out
    if ! diff -q mine.out brute.out > /dev/null; then
        echo "第 $i 組不一樣"
        cat test.in
        exit 1
    fi
done
echo "$4 組都一樣"
\end{code}

先 \texttt{chmod +x check.sh} 讓它可以執行，再用 \texttt{./check.sh vote vote\_brute
vote\_gen 1000} 跑一千組。 範例~\ref{ex:vote} 的 \texttt{vote} 會印出
\texttt{1000 組都一樣}。

來看它抓到錯誤的樣子。 有人用全域的計數陣列，每組結束時把用過的格子清掉，可是清的
範圍寫錯了：

\begin{code}[label=code:vote-bug]{cpp}{清錯範圍的計數（錯誤示範）}
int cnt[500001];                          // 全域陣列，一開始全是 0
// ……讀票、找最大值、輸出，和 vote.cpp 相同……
for (int c = 1; c <= n; c++) {            // 應該是 c <= k
    cnt[c] = 0;
}
\end{code}

把 $k$ 寫成 $n$ 是很容易打錯的地方，而且兩個範例都會通過：範例裡每一組清掉的範圍剛
好都蓋住了投過票的主題。 用腳本比對，在本書測試的環境下第 46 組就出現不同：

\begin{code}{text}{check.sh 印出的輸入}
第 46 組不一樣
1 3
3
1 4
4
1 2
2
0 0
\end{code}

暴力解印出 \texttt{3}、\texttt{4}、\texttt{2}，這支程式印出 \texttt{3}、\texttt{3}、
\texttt{2}。 第一組 $n = 1$，只清掉 \texttt{cnt[1]}，主題 $3$ 的那一票留到了第二組，
和主題 $4$ 同票，於是輸出較小的 $3$。 這組輸入只有七行，手算一遍就能看出錯在哪裡。

用暴力解比對時，幾件事要注意：

\begin{itemize}
\item \textbf{暴力解要直接。} 照題目的字面意思枚舉，不要在暴力解裡用同一個「聰明的觀
察」。 兩邊都用了同一個錯的想法，就會一起錯，比對不出來。
\item \textbf{兩邊不同時，先手算。} 有時候錯的是暴力解，或是產生器產生了題目不允許的輸
入。
\item \textbf{比對全部相同不等於證明。} 範圍小的隨機資料抓不到溢位和浮點數誤差，也抓
不到太慢。 這些要靠上一小節的最大輸入，以及下面的工具。
\end{itemize}

迴圈次數的例題也能這樣做：暴力解就是第~\ref{sec:ex-loop-count}~小節照字面跑的那支程式，產生器隨機產生
$n \le 1000$ 的詢問。 不過它的浮點數錯誤只在 $n$ 超過 $2^{53}$ 時才會出現，暴力解根本
跑不到那裡，只能靠手造的輸入（$n = 2^{59} - 1$）抓。 兩種測試方法要一起用。

\subsection{讓編譯器幫忙}

練習時用下面的選項編譯：

\begin{code}{bash}{練習時的編譯選項}
g++ -std=c++17 -O2 -Wall -Wextra -fsanitize=address,undefined vote.cpp -o vote
\end{code}

\texttt{-Wall -Wextra} 讓編譯器對可疑的寫法提出警告，例如比較有號和無號的整數、變數
沒有初始化就使用。 \texttt{-fsanitize=address,undefined} 會在執行時檢查陣列越界和有
號整數溢位，出事時印出是哪一行（銜接書第~10~章 10.1 節有示範）。 第~\ref{sec:overflow}~
節的溢位程式這樣編譯後執行，會印出 \texttt{runtime error: signed integer overflow:
100000 * 100000 cannot be represented in type 'int'} 這樣的訊息，前面還有檔名和行號。
加了這些選項，程式會慢好幾倍，所以只在測試時用；量時間時拿掉。

量時間用 \texttt{time} 指令，例如 \texttt{time ./vote < big.in > /dev/null}，看
\texttt{real} 那一行。 \texttt{big.in} 可以用另一個產生器做出來：把
\texttt{vote\_gen} 的範圍改成題目的上限，就是一個最大輸入的產生器。

\begin{apcsnote}
實作題由系統用測試資料評分，考場上看不到這些測試資料，只有題目的範例。 練習時就養
成習慣：寫完先跑範例，再手造幾筆邊界輸入；時間夠的話，寫一個暴力解在小範圍比對。 考
試時不一定有時間寫完整的腳本，但手造最小、最大、剛好相等這三種輸入，幾分鐘就做得完。
\end{apcsnote}

\section{本章常見錯誤}
\label{sec:prep-pitfalls}

前面各節的錯誤，歸納起來是五類。 寫每一題時，交出去之前逐項檢查一次。

\begin{center}
\begin{tabular}{p{0.2\linewidth}p{0.48\linewidth}p{0.2\linewidth}}
\toprule
錯誤 & 例子 & 見 \\
\midrule
估錯工作量 & 看到兩層迴圈就說 $O(n^2)$ 或 $n \log n$；只看 $n$ 沒看詢問數 $q$；每組都清空整個大陣列 &
第~\ref{sec:ex-loop-count}、\ref{sec:from-limits}、\ref{sec:ex-vote}~節 \\
溢位 & \texttt{int} 乘完才轉 \texttt{long long}；\texttt{1 << b}；比大小時直接相乘 &
第~\ref{sec:overflow}~節 \\
浮點數誤差 & 用 \texttt{sqrt}、\texttt{cbrt}、\texttt{log2} 求整數答案，沒有修正 &
第~\ref{sec:ex-loop-count}、\ref{sec:ex-max-size}~節 \\
輸入輸出 & 沒加速、用 \texttt{endl}；\texttt{>>} 後面直接 \texttt{getline}；結束標記判斷錯；多印空白 &
第~\ref{sec:io}~節 \\
只用範例測試 & 清錯範圍、浮點數誤差、太慢，範例都會通過 &
第~\ref{sec:testing}~節 \\
\bottomrule
\end{tabular}
\end{center}

\section{本章練習}
\label{sec:ch1-exercises}

\begin{drills}[label=drills:ch1]
\item 下面兩段程式，$n = 10$ 時 \texttt{cnt++} 各執行幾次？ 一般的 $n$ 時，複雜度是多少？
\begin{code}{cpp}{練習 1}
// (a)
for (int i = 1; i <= n; i++)
    for (int j = i; j <= n; j += i)
        cnt++;
// (b)
for (int i = n; i >= 1; i /= 2)
    for (int j = 1; j <= n; j++)
        cnt++;
\end{code}

\item 程式甲處理 $n$ 筆資料要 $2n^2$ 步，程式乙要 $200n + 5000$ 步。
\begin{enumerate}
\item $n = 50$ 和 $n = 1000$ 時，哪一支比較快？
\item 找出最小的 $n$，使得從這個 $n$ 開始，程式乙一直都嚴格比程式甲快。
\end{enumerate}

\item 下面每一題，判斷想用的做法來不來得及（用每秒 $10^8$ 次估），並說明理由。
\begin{enumerate}
\item $n \le 2 \times 10^5$ 個數、$q \le 2 \times 10^5$ 個詢問，每個詢問把陣列從頭掃
一遍。
\item 最多 $10^4$ 組資料，每組 $n \le 1000$，每組用 $O(n^2)$ 的做法。
\item 同 (b)，但題目另外保證所有組的 $n$ 加起來不超過 $10^5$。
\item $n \le 20$，枚舉所有子集合，每個子集合花 $O(n)$ 檢查。
\end{enumerate}

\item 下面的程式片段中，$n \le 10^5$，$0 \le$ \texttt{a[i]} $\le 10^6$。 哪幾行可能
溢位？ 怎麼改？
\begin{code}{cpp}{練習 4}
long long sum = 0, sq = 0;
for (int i = 0; i < n; i++) {
    sum += a[i];                       // (1)
    sq += a[i] * a[i];                 // (2)
}
int avg = sum / n;                     // (3)
long long total = n * (n + 1) / 2;     // (4)
\end{code}

\item 為第~\ref{sec:ex-loop-count}~小節的例題設計五筆手造的輸入，每一筆寫出它要檢查的
是哪一種錯誤，以及正確的輸出。
\end{drills}

\section{實作練習}
\label{sec:ch1-practice}

下面是 Aki-Judge 上和本章相關的題目。 難度分四階：暖身是直接套用本章的寫法；基本要
自己想一下怎麼估計或怎麼讀輸入；實戰的題敘比較長，要先把題目轉換成本章的形狀；挑戰
需要再多一步觀察。

\begin{center}
\begin{tabular}{llc}
\toprule
小節 & 題目 & 難度 \\
\midrule
估計次數與複雜度 & \akiref{stage-runs-under-the-entry-cap}{依入口上限計算的連續關卡} & 基本 \\
 & \akiref{items-ordered-exactly-once}{只出現一次的值} & 實戰 \\
 & \akiref{rectangles-with-an-even-number-of-stamps}{章數為偶數的矩形} & 實戰 \\
 & \akiref{partner-pairs-at-two-scales}{兩種規模的搭檔配對} & 實戰 \\
 & \akiref{faster-from-here-on}{從此以後都比較快} & 挑戰 \\
\midrule
整數運算與溢位 & \akiref{sheets-for-a-saddle-stitched-booklet}{騎馬釘小冊子的紙張數} & 暖身 \\
 & \akiref{position-on-a-circular-track}{環形跑道上的位置} & 暖身 \\
 & \akiref{splitting-the-class-trip-bill}{班遊費用分攤} & 基本 \\
 & \akiref{download-progress-bar}{下載進度條} & 基本 \\
 & \akiref{stairs-or-the-express-elevator}{樓梯還是快速電梯} & 實戰 \\
 & \akiref{per-purchase-vs-monthly-loyalty-points}{逐筆給點與月結給點的差距} & 挑戰 \\
\midrule
輸入與輸出 & \akiref{jump-rope-record-breaks}{跳繩計數器的破紀錄次數} & 暖身 \\
 & \akiref{ballots-of-unknown-count}{不知道有幾行的投票紀錄} & 暖身 \\
 & \akiref{steps-walked-during-each-break}{每天各節下課的步數} & 暖身 \\
 & \akiref{drink-shop-sales-per-shift}{飲料店每班的營業額} & 暖身 \\
 & \akiref{study-group-sign-in-sheets}{讀書會簽到表} & 基本 \\
 & \akiref{study-room-seat-log}{自習室的座位登記} & 實戰 \\
\bottomrule
\end{tabular}
\end{center}

本章的三題例題也在 Aki-Judge 上，讀完各節之後可以先自己寫一次再對照。

下面這組是 STL 的複習題，用到的是銜接書第~3、8、9~章的 \texttt{vector}、
\texttt{pair}、\texttt{sort} 與比較函式、\texttt{set}／\texttt{map}。 從第~2~章起，
這些容器會直接拿來用，不再說明。 寫起來卡住的話，先回頭翻銜接書對應的那一章。

\begin{center}
\begin{tabular}{llc}
\toprule
小節 & 題目 & 難度 \\
\midrule
排序與比較函式 & \akiref{sorting-the-camp-roster}{營隊報名名冊排序} & 暖身 \\
 & \akiref{graduation-photo-rows}{畢業照的分排站位} & 暖身 \\
 & \akiref{league-ranks-by-win-rate}{校際聯賽勝率名次} & 基本 \\
 & \akiref{sorting-file-names-the-human-way}{照人的直覺排檔名} & 基本 \\
 & \akiref{clicking-column-headers}{試算表連點欄位標題} & 實戰 \\
 & \akiref{winter-break-reading-marathon-standings}{寒假閱讀馬拉松結算} & 實戰 \\
 & \akiref{order-of-the-lighthouse-sweep}{燈塔光束掃過的順序} & 挑戰 \\
\midrule
\texttt{set} 與 \texttt{map} & \akiref{hallway-light-switch-log}{走廊燈的開關紀錄} & 暖身 \\
 & \akiref{the-calculators-variable-table}{小算盤的變數表} & 暖身 \\
 & \akiref{class-bonus-and-personal-bonus}{全班加分與個人加分} & 基本 \\
 & \akiref{cone-squares-on-the-field}{操場上的方格標記} & 基本 \\
 & \akiref{transposed-or-mirrored}{移調與倒影} & 實戰 \\
 & \akiref{matching-prefixes-of-two-attendance-lists}{兩份點名單的前綴配對} & 實戰 \\
 & \akiref{left-and-right-socks-in-a-row}{左右腳都配成雙的區段} & 挑戰 \\
\bottomrule
\end{tabular}
\end{center}

\section{讀完你應該會}
\label{sec:ch1-checklist}

\begin{itemize}
\item 給一段內層範圍跟著外層變、或變數每次乘除一個數的迴圈，能列表算出它的精確執行次
數，並寫出它的複雜度。
\item 看到題目的限制，能把 $n$、$q$、組數都代進自己的做法，判斷來不來得及；題目限制
「總和」時，能讓每一組的工作量跟那一組的大小成正比。
\item 寫程式前能估出答案和中間值的最大值，在需要的地方用 \texttt{long long}、
\texttt{1LL}，並用除法比較可能溢位的乘積。
\item 會用 \texttt{while (cin >> x)} 讀到檔案結束、用 \texttt{getline} 加
\texttt{istringstream} 逐行拆資料，並正確處理多組資料的結束標記。
\item 能為一題寫出暴力解、隨機產生器和比對腳本，找出一組讓兩支程式輸出不同的小輸入。
\end{itemize}

\textbf{延伸閱讀}（書末會整理完整的參考資料）

\begin{itemize}
\item Antti Laaksonen，\textit{Competitive Programmer's Handbook}，第 1、2 章：輸入輸
出、整數型別和時間複雜度的簡短整理，並用最大子陣列和示範一個做法怎麼一步步變快。
\item Steven Halim 等，\textit{Competitive Programming 4}，Book 1 的 1.3 節：怎麼從
限制估計複雜度，以及測試時該準備哪些輸入。
\item Steven S. Skiena，\textit{The Algorithm Design Manual}，第 3 版，第 2 章：
big-O 的定義和常見的成長速度，寫得比本章嚴謹。
\end{itemize}

\addcontentsline{toc}{section}{練習解答}
\section*{練習解答}

\begin{enumerate}
\item
\begin{enumerate}
\item $27$ 次。 外層是 \texttt{i} 時，\texttt{j} 是 $i, 2i, 3i, \ldots$，共
$\lfloor n / i \rfloor$ 次，$n = 10$ 時是 $10 + 5 + 3 + 2 + 2 + 1 + 1 + 1 + 1 + 1
= 27$。 一般的 $n$ 是 $\lfloor n/1 \rfloor + \lfloor n/2 \rfloor + \cdots +
\lfloor n/n \rfloor$，大約是 $n$ 乘上 $1 + 1/2 + \cdots + 1/n$，而這個和大約是
$\ln n$，所以是 $O(n \log n)$。
\item $40$ 次。 外層 \texttt{i} 是 $10, 5, 2, 1$，共 $4$ 輪，每輪內層 $10$ 次。 一般
的 $n$ 是 $n \times (\lfloor \log_2 n \rfloor + 1)$，$O(n \log n)$。
\end{enumerate}

\item
\begin{enumerate}
\item $n = 50$ 時甲要 $5000$ 步，乙要 $15000$ 步，甲比較快。 $n = 1000$ 時甲要
$2 \times 10^6$ 步，乙要 $205000$ 步，乙比較快。
\item $121$。 $n = 120$ 時甲 $28800$、乙 $29000$，甲還是比較快；$n = 121$ 時甲
$29282$、乙 $29200$，乙比較快。 兩者的差 $2n^2 - 200n - 5000$ 在 $n \ge 50$ 之後只
會越來越大，所以從 $121$ 起乙一直比較快。
\end{enumerate}

\item
\begin{enumerate}
\item 不行。 $n \times q = 4 \times 10^{10}$。 要換成第~5~章的前綴和這類先處理好、
每個詢問 $O(1)$ 或 $O(\log n)$ 的做法。
\item 不行。 每組最多 $10^6$ 步，$10^4$ 組是 $10^{10}$。
\item 可以。 每組 $n_i^2 \le n_i \times 1000$，加起來不超過
$1000 \times \sum n_i = 1000 \times 10^5 = 10^8$。 總和的限制讓「很多組都很大」的情況
不會發生。
\item 可以。 $2^{20} \times 20 \approx 2 \times 10^7$。
\end{enumerate}

\item 第 (2) 行和第 (4) 行。
\begin{itemize}
\item (2) \texttt{a[i] * a[i]} 用 \texttt{int} 乘，最大 $10^{12}$，溢位。 改成
\texttt{(long long)a[i] * a[i]}。 加起來最多 $10^{17}$，\texttt{sq} 存得下。
\item (1) 和 (3) 沒問題：\texttt{sum} 是 \texttt{long long}，最多 $10^{11}$；平均值不超過 $10^6$。
\item (4) \texttt{n * (n + 1)} 用 \texttt{int} 乘，$n = 10^5$ 時大約 $10^{10}$，溢位。
改成 \texttt{(long long)n * (n + 1) / 2}。
\end{itemize}

\item 一種答法（輸入的第一行 $q$ 省略）：
\begin{itemize}
\item \texttt{3 1}：$n$ 最小，片段 3 的 \texttt{j} 從 $0$ 開始，答案 $2$。 漏掉
$+1$ 會印 $1$。
\item \texttt{1 1}：片段 1 只有一格，答案 $1$。
\item \texttt{2 7} 和 \texttt{2 8}：$n$ 剛好在 $2$ 的次方前後，答案 $7$ 和 $15$，檢查
迴圈條件的 $\le$。
\item 片段 1 取最大的 $n$（\texttt{1} 後面接 $10^{18}$）：檢查時間和溢位，答案是
\texttt{1999999999999999976}。
\item \texttt{2 576460752303423487}：$n = 2^{59} - 1$，檢查有沒有用浮點數算輪數，答案
是 $576460752303423487$。
\end{itemize}
\end{enumerate}
